\section{Orphaned statements}
\label{sec:orphaned-statements}

The blueprint's spine --- every statement that flows into the challenge API
at \ref{thm:challenge-api} --- does not require these results. They are
classical facts that the formalisation infrastructure happens to be able
to support, but they are not on the path to discharging the
\code{Jacobian/Challenge.lean} obligations. They live here so that (a)
the work that has gone into them is preserved, and (b) the connectivity
auditor (\code{scripts/blueprint\_graph\_audit.py}) can ignore them by
recognising the \code{sec:orphaned-statements} marker.

If a statement in this section is later needed by the spine, move it back
into the appropriate numbered section and add the \code{\textbackslash uses}
edge that brings it in. The goal is to prove exactly what the challenge
API requires; nothing here counts toward that goal.

\subsection{Riemann's inequality}

\begin{theorem}[Riemann's inequality]
\label{thm:riemann-inequality}
\uses{thm:euler-char-line-bundle}
For any line bundle \(\mathcal{L}\) on a compact Riemann surface of
genus \(g\),
\[
  h^{0}(\mathcal{L}) \ge \deg(\mathcal{L}) + 1 - g .
\]
The ``weaker half'' of Riemann-Roch (1857), preceding Roch's identification
of the \(h^{1}\) correction term: the inequality drops the non-negative
\(h^{1}(\mathcal{L})\) summand from the Euler characteristic identity.
The challenge API uses the full Riemann-Roch identity
(\ref{input:riemann-roch}), which already implies this inequality, so
this weaker statement is redundant on the spine.
\lean{JacobianChallenge.HolomorphicForms.riemann_inequality_h0}
\leanok
\end{theorem}
\begin{proof}\leanok\end{proof}

\subsection{Riemann's theorem on the theta divisor}

\begin{theorem}[Riemann's theorem on the theta divisor]
\label{thm:riemann-theorem-theta-divisor}
\uses{def:analytic-jacobian,input:riemann-roch}
For \(g\ge 1\), a point \(z\in\Jac(X)\) lies on the theta divisor
\(\Theta\) if and only if \(z = \kappa + \AJ_{g-1}(s)\) for some
\(s\in\operatorname{Sym}^{g-1}(X)\), where \(\kappa\) is the Riemann
constant and \(\AJ_{g-1}\) is the Abel-Jacobi map on symmetric
\((g-1)\)-fold products.
\depends{Mumford, Tata Lectures on Theta II, \S3. Mathlib hooks: theta
function on Siegel upper half space, Riemann-Roch for line bundles
attached to divisors on \(X\).}
The challenge API does not mention the theta divisor, so this theorem
is off-spine.
\lean{JacobianChallenge.Blueprint.AbelExistence.RiemannTheoremThetaDivisor.riemann_theorem_theta_divisor}
\leanok
\end{theorem}
\begin{proof}\leanok\end{proof}
